\begin{lemma}[Better quasi-order implies well quasi-order]
For every type $Q$ and relation $r$ on $Q$,
\[
\mathsf{Bqo}(r)\longrightarrow \mathsf{WellQuasiOrdered}(r).
\]
No reflexivity, transitivity, or other order hypothesis is required for this
implication.
\end{lemma}
\begin{proof}
Assume $\mathsf{Bqo}(r)$ and suppose, toward a contradiction, that
$r$ is not well-quasi-ordering.  By the sequence definition of
$\mathsf{WellQuasiOrdered}$ from the shared environment, there is a sequence
$a:\mathbb N\to Q$ such that for all $m<n$,
\[
\neg r(a(m),a(n)).
\]
Define a multi-sequence $h:\mathsf{IncSeq}\to Q$ by
\[
h(x)=a(x(0)).
\]
This has the required type because \noderef{MultiSequence} identifies a
multi-sequence with a map out of \noderef{IncSeq}'s increasing
enumerations.

We first check local constancy.  For any $x$, take the one-entry prefix
length $1$.  If $y$ agrees with $x$ on this prefix, then
\noderef{PrefixAgree} gives $y(0)=x(0)$, and hence
$h(y)=a(y(0))=a(x(0))=h(x)$.  Thus the definition of
\noderef{LocallyConstant} is satisfied.

We next check badness.  For any increasing enumeration $x$, strict
monotonicity gives $x(0)<x(1)$.  The assumed badness of $a$ therefore yields
\[
\neg r(a(x(0)),a(x(1))).
\]
By the definition of \noderef{ShiftMap}, the shifted enumeration has first
entry $x(1)$, so this is exactly
\[
\neg r(h(x),h(\mathsf{ShiftMap}(x))).
\]
Consequently \noderef{BadMultiSequence} says that $h$ is a bad
multi-sequence.  Its local constancy and badness contradict the defining
exclusion in \noderef{Bqo}.  Hence $r$ is well-quasi-ordered.
\end{proof}
